% LỜI TỰA --- Quyển II
\chapter*{Lời Tựa}
\addcontentsline{toc}{chapter}{Lời Tựa}
\markboth{Lời Tựa}{Lời Tựa}

\begin{truyen}{Lời Tựa}{Tales of Character and Human Values}
\chuhoa{G}{iáo viên bắt đầu bộ sách này với một câu hỏi đơn giản: làm thế nào để mỗi buổi học có thể để lại thứ gì đó đáng nhớ hơn là bài thi?}
\textit{(A teacher began this series with a simple question: how can each class leave behind something more memorable than a test?)}

Đạo làm người không học một lần rồi thôi~--- nó cần được nhắc đi nhắc lại suốt đời.
\textit{(The way of being human is not learned once and forgotten~--- it needs to be revisited throughout life.)}

Quyển hai của bộ sách tiếp tục hành trình đó. Những câu chuyện ở đây xoay quanh đạo làm con, làm bạn, làm trò, làm thầy, làm người trong đời. Mỗi bài học được đặt trong một tình huống quen thuộc để người đọc dễ liên hệ với cuộc sống thực của chính mình.
\textit{(The second volume continues that journey. The stories here revolve around being a good child, friend, student, teacher, and human being. Each lesson is placed in a familiar situation so readers can easily connect it to their own real lives.)}

Bộ sách <<Truyện Tích Cảm Hứng Song Ngữ>> gồm 24 quyển, mỗi quyển một chủ đề riêng. Quyển II này mang chủ đề: \textbf{đạo làm con, làm bạn, làm trò, làm thầy}.
\textit{(The series ``Inspiring Bilingual Tales'' contains 24 volumes, each with its own theme. Volume II focuses on: \textbf{đạo làm con, làm bạn, làm trò, làm thầy}.})
\end{truyen}

\begin{baihoc}
Nhân cách không hình thành trong một ngày~--- nhưng mỗi ngày đều có thể là ngày bắt đầu.
\textit{(Character is not formed in a day~--- but every day can be the day it begins.)}
\end{baihoc}

\begin{ghinhoanh}
\textbf{Short English to remember:}

Good character is built one day at a time.

The way we treat others reflects who we are.
\end{ghinhoanh}

\begin{tuhocnhanh}
1. Trong các vai trò~--- con, bạn, trò, thầy~--- em đang làm tốt nhất ở vai nào? \textit{(Among the roles~--- child, friend, student, teacher~--- which do you play best?)}

2. Câu chuyện nào trong quyển này khiến em nhìn lại cách ứng xử của mình? \textit{(Which story here made you reconsider how you treat others?)}

3. Điều gì trong đạo làm người em muốn rèn luyện thêm? \textit{(Which aspect of good character do you want to improve?)}
\end{tuhocnhanh}

\begin{center}
\textit{\color{truyengold}``Mỗi câu chuyện là một tấm gương~--- hãy nhìn vào để thấy mình.''}
\end{center}

\begin{flushright}
\textbf{Tôi Nguyễn Văn Sang}\\
\textit{GV THPT Nguyễn Hữu Cảnh - TP HCM}
\end{flushright}

\ngancach
